\documentclass[11pt]{article}

\usepackage[margin=1.1in]{geometry}
\usepackage{amsmath, amssymb, mathtools}
\usepackage{bm}
\usepackage[colorlinks=true, linkcolor=blue, citecolor=blue]{hyperref}

\newcommand{\ii}{\mathrm{i}}
\renewcommand{\Re}{\operatorname{Re}}
\renewcommand{\Im}{\operatorname{Im}}
\newcommand{\hc}{\mathrm{h.c.}}
\newcommand{\code}[1]{\texttt{#1}}

\title{Frame conversions for the CNOT3 convergence experiment:\\
rotating frame with and without the RWA, and the laboratory frame}
\author{}
\date{\today}

\begin{document}
\maketitle

\section{Overview}

The CNOT3 problem from the High-Order Hermite Optimization (HOHO) paper is
posed in the rotating frame with the rotating wave approximation (RWA)
applied, and its controls were optimized in that setting.  The convergence
experiment additionally solves the \emph{same physical problem} in two other
forms:
\begin{description}
    \item[\code{rwa}] rotating frame, RWA applied (the original setting);
    \item[\code{norwa}] rotating frame, counter-rotating control terms
        retained (no RWA);
    \item[\code{lab}] laboratory frame.
\end{description}
The \code{norwa} and \code{lab} dynamics are \emph{exactly} equivalent ---
related by a known diagonal unitary --- while \code{rwa} differs from both by
the physical error of the rotating wave approximation.  This note derives the
conversions and records how they are implemented as transformations of
QuantumGateDesign \code{CarrierControl} objects
(\code{drop\_rwa} and \code{to\_lab\_frame} in
\code{src/QuantumGateDesign\_interface.jl}).

\section{Setup and notation}

The system consists of $Q = 3$ subsystems (storage cavity and two transmon
qubits) with lowering operators $a_k$ and angular transition frequencies
$\omega_k = 2\pi f_k$.  In the laboratory frame,
\begin{equation} \label{eq:lab_hamiltonian}
    H^{\mathrm{lab}}(t)
    = \underbrace{W + H_{\mathrm{Kerr}}}_{\text{drift}}
    + \sum_{k=1}^{Q} f_k(t)\,\bigl(a_k + a_k^\dagger\bigr),
    \qquad
    W \coloneqq \sum_{k=1}^{Q} \omega_k\, a_k^\dagger a_k,
\end{equation}
where $H_{\mathrm{Kerr}}$ collects the (diagonal, number-conserving) self- and
cross-Kerr terms and each $f_k(t)$ is a \emph{real} drive.  Following the
Juqbox/HOHO convention, the drive is written in terms of a complex envelope
$\alpha_k(t) = p_k(t) + \ii q_k(t)$ modulating the carrier $e^{\ii\omega_k t}$:
\begin{equation} \label{eq:lab_pulse}
    f_k(t)
    = 2\,\Re\!\left[\alpha_k(t)\, e^{\ii \omega_k t}\right]
    = \alpha_k(t)\, e^{\ii \omega_k t} + \overline{\alpha_k(t)}\, e^{-\ii \omega_k t}.
\end{equation}
In the CNOT3 problem each envelope is itself a sum of slowly rotating
carriers,
\begin{equation} \label{eq:carrier_control}
    \alpha_k(t) = \sum_{f} E_{k,f}(t)\, e^{\ii \Omega_{k,f} t},
    \qquad
    E_{k,f}(t) = P_{k,f}(t) + \ii\, Q_{k,f}(t),
\end{equation}
with B-spline envelopes $P_{k,f}, Q_{k,f}$ and carrier frequencies
$\Omega_{k,f}$ at the (negated) Kerr shifts.  Equation
\eqref{eq:carrier_control} is precisely what a QuantumGateDesign
\code{CarrierControl} evaluates: it returns
$p_k = \Re \alpha_k$ and $q_k = \Im \alpha_k$, and stores each $E_{k,f}$ as a
coefficient slice $[P_{k,f};\, Q_{k,f}]$.  Two consequences used below:
adding a carrier means appending a frequency and a slice, and conjugating an
envelope, $\overline{E} = P - \ii Q$, means copying a slice and negating its
$Q$-half.

\section{Rotating frame and the RWA}

Transform with the diagonal unitary $R(t) = e^{\ii W t}$, so
$\bm\psi_R = R\,\bm\psi^{\mathrm{lab}}$ and
\begin{equation} \label{eq:frame_change}
    H_R(t) = R H^{\mathrm{lab}} R^\dagger - W,
    \qquad
    R\, a_k R^\dagger = a_k\, e^{-\ii \omega_k t}.
\end{equation}
The $W$ in the lab drift cancels and $H_{\mathrm{Kerr}}$ commutes with $R$, so
the rotating-frame drift is $H_{\mathrm{Kerr}}$ alone.  The control term
becomes, using \eqref{eq:lab_pulse},
\begin{align}
    f_k(t)\bigl(a_k e^{-\ii\omega_k t} + a_k^\dagger e^{\ii\omega_k t}\bigr)
    &= \bigl(\alpha_k e^{\ii\omega_k t} + \overline{\alpha_k} e^{-\ii\omega_k t}\bigr)
       \bigl(a_k e^{-\ii\omega_k t} + a_k^\dagger e^{\ii\omega_k t}\bigr) \nonumber \\
    &= \underbrace{\alpha_k\, a_k + \overline{\alpha_k}\, a_k^\dagger}_{\text{co-rotating}}
     \;+\; \underbrace{\overline{\alpha_k}\, e^{-2\ii\omega_k t}\, a_k
         + \alpha_k\, e^{2\ii\omega_k t}\, a_k^\dagger}_{\text{counter-rotating}}.
     \label{eq:rotating_control}
\end{align}
The RWA discards the counter-rotating terms, which oscillate at $2\omega_k$.
What remains is the \code{rwa} problem as QuantumGateDesign poses it:
\begin{equation} \label{eq:rwa_hamiltonian}
    H^{\mathrm{rwa}}(t)
    = H_{\mathrm{Kerr}}
    + \sum_k \Bigl[\, p_k(t)\,\bigl(a_k + a_k^\dagger\bigr)
        + q_k(t)\, \ii\bigl(a_k - a_k^\dagger\bigr) \Bigr],
\end{equation}
since $\alpha_k a_k + \overline{\alpha_k} a_k^\dagger
= p_k (a_k + a_k^\dagger) + q_k\, \ii (a_k - a_k^\dagger)$.  The factor $2$ in
\eqref{eq:lab_pulse} is fixed by this requirement: with it, the optimized
control vector of the \code{rwa} problem describes the same physical pulse in
all three frames.

\section{\code{norwa}: keeping the counter-rotating terms}

Keeping all of \eqref{eq:rotating_control}, the control term retains the
structure $\gamma_k a_k + \overline{\gamma_k} a_k^\dagger$ with a modified
complex envelope
\begin{equation} \label{eq:gamma}
    \boxed{\;
    \gamma_k(t) = \alpha_k(t) + \overline{\alpha_k(t)}\, e^{-2\ii\omega_k t}
    \;}
\end{equation}
replacing $\alpha_k$; the drift is unchanged.  Substituting
\eqref{eq:carrier_control},
\begin{equation}
    \gamma_k(t)
    = \sum_f E_{k,f}(t)\, e^{\ii \Omega_{k,f} t}
    + \sum_f \overline{E_{k,f}(t)}\; e^{-\ii (2\omega_k + \Omega_{k,f}) t},
\end{equation}
which is again of the form \eqref{eq:carrier_control}.  \code{drop\_rwa}
therefore rebuilds each control with the doubled carrier list
\begin{equation}
    \bigl\{\Omega_{k,f}\bigr\}_f \;\cup\;
    \bigl\{-(2\omega_k + \Omega_{k,f})\bigr\}_f ,
\end{equation}
where the appended carriers carry the conjugated coefficient slices
$[P_{k,f};\, -Q_{k,f}]$.  Nothing else changes: same operators, same drift,
same solvers.

\section{\code{to\_lab\_frame}: undoing the rotation}

In the laboratory frame \eqref{eq:lab_hamiltonian}, the drift regains $W$
(implemented by rebuilding the problem with zero rotation frequencies, so the
transition frequencies appear on the drift diagonal) and the control couples
\emph{only} through the symmetric operator $a_k + a_k^\dagger$ with the real
pulse $f_k$; there is no antisymmetric part, so the problem's antisymmetric
control operators are set to zero.  It remains to express $f_k$ as the
$p$-part of a \code{CarrierControl}.  From \eqref{eq:lab_pulse} and
\eqref{eq:carrier_control},
\begin{equation}
    f_k(t)
    = 2\,\Re\!\left[\sum_f E_{k,f}(t)\, e^{\ii(\omega_k + \Omega_{k,f}) t}\right]
    = \Re\!\left[\sum_f 2E_{k,f}(t)\, e^{\ii(\omega_k + \Omega_{k,f}) t}\right],
\end{equation}
so \code{to\_lab\_frame} shifts every carrier to $\omega_k + \Omega_{k,f}$ and
doubles the coefficients.  The resulting control's $q$-part,
$\Im[\sum_f 2E_{k,f} e^{\ii(\omega_k+\Omega_{k,f})t}]$, is \emph{not} zero ---
it is a by-product of reusing the $p + \ii q$ machinery --- but it is
harmless because it multiplies the zeroed antisymmetric operator.

\section{Equivalence and verification}

Because \eqref{eq:rotating_control} is an exact rewriting of the lab-frame
control, the \code{norwa} and \code{lab} solutions satisfy, for all $t$,
\begin{equation} \label{eq:equivalence}
    \bm\psi^{\mathrm{norwa}}(t) = e^{\ii W t}\, \bm\psi^{\mathrm{lab}}(t),
\end{equation}
with $W$ recoverable as the difference of the two drift diagonals (the Kerr
terms cancel).  The script \code{scripts/cnot3/verify\_frames.jl} checks:
\begin{enumerate}
    \item the envelope identities \eqref{eq:gamma} and \eqref{eq:lab_pulse}
        pointwise, and control derivatives against finite differences;
    \item the solution identity \eqref{eq:equivalence} at the final time on a
        reduced system, where it holds to the solver error
        ($\sim 10^{-8}$);
    \item that the \code{rwa} solution differs from \code{norwa} by a small
        but nonzero amount --- the physical RWA error, $\sim 10^{-3}$ on the
        reduced system --- confirming that the counter-rotating terms are
        present and have the expected (small) physical effect.
\end{enumerate}

Numerically, the frames stress the integrators differently.  In \code{rwa}
all time dependence is slow.  In \code{norwa} the drift is unchanged but the
controls acquire carriers at $2\omega_k$; in \code{lab} the fast oscillation
appears both in the drift diagonal (handled exactly by the Filon ansatz
$u(t) = g(t)e^{\ii\omega t}$) and in the control carriers at
$\omega_k + \Omega_{k,f}$.  Methods that sample the controls through
polynomial interpolation (Hermite, plain Filon) must resolve those carriers
with the timestep, while the controlled-Filon method factors them into the
quadrature and is insensitive to them.
\end{document}
